\section{Conclusion}
\label{sec:conclusion}

Organizing a decade of COF photocatalysis along a single causal chain --
linkage $\rightarrow$ conjugation $\rightarrow$ band structure $\rightarrow$
exciton and charge separation $\rightarrow$ interface $\rightarrow$ HER --
turns a sprawling literature into a set of testable design rules. Linkage
chemistry no longer forces a stability--conjugation trade-off: vinylene and
sp$^2$-carbon frameworks deliver both, while $\beta$-ketoenamine and imine
platforms remain the workhorses for systematic study~\cite{jin20192d,
cheng2023fully,xue2023research}. Donor--acceptor dipole engineering,
verified by exciton-binding measurements rather than nominal composition,
and linkage protonation are the most portable band-structure levers
~\cite{qian2023computation,li2022covalent,yang2021protonated}. Mechanistic
claims increasingly rest on multi-probe convergence with matched controls
~\cite{ghosh2020identification,chu2024controlled}, and the strongest
performance reports anchor rates to AQY($\lambda$) with explicit protocol
variables~\cite{qiu2026atomically,zhao2024nanoscale}. What stands between
laboratory quantum efficiency and solar fuels is no longer molecular
designability but verification infrastructure: portable metrics,
donor-free operation, thousand-hour stability, scalable synthesis, and
reactor-level photon management~\cite{sun2025design,wu2025recent}. The
pre-registered TpPa-1 route of Section~\ref{sec:idea} is offered as a
template for how the field can convert its accumulated design rules into
gated, falsifiable experiments -- and how negative results, bounded by
matched controls, can constrain the causal chain as productively as
records~\cite{mo2020alkene,wang2026bandgap}.
